\documentclass{article}
\usepackage[utf8]{inputenc}

% ECS 240 HW1 Problem Set

\usepackage{amsmath}
\usepackage{amssymb}
\usepackage{amsthm}
\usepackage{enumerate}

\usepackage[margin=1.5in]{geometry}

\usepackage{xcolor}
\usepackage{hyperref}
\hypersetup{
    colorlinks,
    linkcolor={red!50!black},
    citecolor={blue!50!black},
    urlcolor={blue!80!black}
}

\usepackage{booktabs,tabularx}

% Keep words intact instead of hyphenating them at the end of a line.
\hyphenpenalty=10000
\exhyphenpenalty=10000
\emergencystretch=3em

% \newcommand{\visiblehref}[2]{\href{#1}{#2}\footnote{\url{#1}}}

\begin{document}

\title{HW1 Problem Set}
\author{ECS 240}
\date{Due: Friday, October 23, 2026}

\maketitle

\hrule

\bigskip

\noindent
This problem set contains 5 problems.
Problems 1, 3, and 5 are coding problems,
and problems 2 and 4 are pure theory problems.
Problems 2, 3, and 5 build on Problem 1,
while Problem 4 is standalone.
Most problems have multiple parts.
There is a template in \texttt{src/} which you can modify for the programming problems.
Solutions to the theory problems should be written in LaTeX.

\bigskip

\hrule

\bigskip

\begin{enumerate}

%%%%%%%%%%%%%%%%%%%%%%%%%
%%%%%   PROBLEM 1   %%%%%
%%%%%%%%%%%%%%%%%%%%%%%%%

\item
``Brainfuck'' (we shall call it BF) is a toy programming language created by a Swiss computer programmer, Urban Müller, in 1993.
It's known as one of the simplest ``esoteric'' programming languages, which means it is intended to be excessively simple and/or convoluted to write programs in, but it can be used to express arbitrary computation.
It is also related to several other esoteric programming languages.

There are eight commands in BF. The following table is reproduced from Wikipedia:\footnote{
    \url{https://en.wikipedia.org/wiki/Brainfuck}
}

\begin{table}[ht]
\centering
\begin{tabularx}{\linewidth}{@{}cX@{}}
\toprule
Character & Semantics \\
\midrule
\texttt{>} & Increment the data pointer by one (to point to the next cell to the right). \\
\texttt{<} & Decrement the data pointer by one (to point to the next cell to the left). Undefined if at 0. \\
\texttt{+} & Increment the byte at the data pointer by one modulo 256. \\
\texttt{-} & Decrement the byte at the data pointer by one modulo 256. \\
\texttt{.} & Output the byte at the data pointer. \\
\texttt{,} & Accept one byte of input, storing its value in the byte at the data pointer. \\
\texttt{[} & If the byte at the data pointer is zero, then instead of moving the instruction pointer forward to the next command, jump it \emph{forward} to the command after the \emph{matching} \texttt{]} command. \\
\texttt{]} & If the byte at the data pointer is nonzero, then instead of moving the instruction pointer forward to the next command, jump it \emph{back} to the command after the \emph{matching} \texttt{[} command. \\
\bottomrule
\end{tabularx}
\end{table}

The syntax of BF can be any sequence of command characters, but with one important restriction:
the left and right brackets (\texttt{[} and \texttt{]}) have to be well-matched. For example:
$$
\texttt{[><+-[.]][+]}
$$
is a valid BF program, but
$$
\texttt{]><+[]}
$$
is not.

\begin{enumerate}[(a)]
\item
Define the syntax for BF programs using a Rust enum, \texttt{enum BFProg}.
You will want to have cases for the different syntax elements above, such as \texttt{IncData} for incrementing the data pointer, \texttt{DecData} for decrementing the data pointer, etc.

\textbf{Hint/notes:} There are multiple ways to define this. However, the important question is how to ensure
the ``left and right brackets are well-matched'' requirement described above.
In Rust - we can model this using one of the recursive cases (in our Rust enums).
You should have only one recursive case of your Rust enum to correspond to the left/right bracket loops,
such as a \texttt{Loop} case.
You will also need a way to combine BF programs in sequence (like \texttt{Sequence} from class.)

\item
Give a way of printing out BF programs by implementing the \texttt{Display} trait for your \texttt{BFProg} data type.
Display in Rust is a way to tell Rust how to print out a datatype as a string/text to the terminal.\footnote{\url{https://doc.rust-lang.org/std/fmt/trait.Display.html}}
\begin{verbatim}
use std::fmt::Display;
impl Display for BFProg {
    fn fmt(&self, f: &mut fmt::Formatter<'_>) -> fmt::Result {
        // ...
    }
}
\end{verbatim}

Your Display trait should print out the BF programs as a sequence of symbols, the way they are given above and in the Wikipedia examples. For example, they should print out like
$$
\texttt{[><+-[.]][+]}
$$

In your \texttt{main()} function, print out a couple of BF programs to see that they are displaying correctly.

\item
Write two unit tests for your syntax which attempt to construct a BF program for the two examples above:
\begin{align*}
&\texttt{[><+-[.]][+]} \\
&\texttt{]><+[]}
\end{align*}

In your first unit test: assert that the first test program prints out correctly.
\begin{verbatim}
assert_eq!(prog.to_string(), "[><+-[.]][+]".to_string());
\end{verbatim}

In your second unit test: you should be unable to construct that exact BF program; you may construct a related program
and assert that it does \emph{not} display as the string above.
\begin{verbatim}
assert!(prog.to_string() != "]><+[]".to_string());
\end{verbatim}

\item
Define a reference interpreter for BF programs.
This should be defined as a \texttt{BFInterpreter} struct/object
with a method \texttt{interpret()},
where the fields of the struct store
the relevant program state,
similar to how we defined a method \texttt{interpret()} on \texttt{VariableStore}
to interpret \texttt{RandProg} programs in class.
The code structure should look like this:
\begin{verbatim}
pub struct BFInterpreter {
    // ... define fields
}
impl BFInterpreter {
    fn new() -> Self {
        // TODO
        unimplemented!()
    }
    fn interpret(&mut self, prog: &BFProg) {
        // TODO
        unimplemented!()
    }
}
\end{verbatim}

\item
Write two unit tests which show that your reference interpreter is working correctly.
The first should be an example BF program which adds two input numbers (like 5 + 6).
The second one can be any example of your choosing.

\begin{verbatim}
#[test]
fn test_bf_interpreter_1() {
    unimplemented!()
}
#[test]
fn test_bf_interpreter_2() {
    unimplemented!()
}
\end{verbatim}
\end{enumerate}

% %%%%%%%%%%%%%%%%%%%%%%%%%
% %%%%%   PROBLEM 2   %%%%%
% %%%%%%%%%%%%%%%%%%%%%%%%%

\item
This problem continues the previous problem.

\begin{enumerate}
\item
In LaTeX, give a formal grammar for the syntax of BF. You may want to review the definition of a
context-free grammar and examples.\footnote{The following may be helpful to see examples of syntax: \url{https://cs.uwaterloo.ca/~s4bendav/files/CS360S22Lec08.pdf} and \url{https://en.wikipedia.org/wiki/Context-free_grammar}}
Your formal grammar should be based on, but not necessarily identical to, your Rust enum from problem 1.

\paragraph{Answer.}
Let the terminal alphabet be
\[
\Sigma = \{\texttt{>},\texttt{<},\texttt{+},\texttt{-},
          \texttt{.},\texttt{,},\texttt{[},\texttt{]}\}.
\]
I use two nonterminals: $P$ for a complete program and $C$ for one
command.  The start symbol is $P$, and the productions are
\[
\begin{aligned}
P &::= \epsilon \mid C\,P,\\
C &::= \texttt{>} \mid \texttt{<} \mid \texttt{+} \mid \texttt{-}
      \mid \texttt{.} \mid \texttt{,} \mid \texttt{[}\,P\,\texttt{]}.
\end{aligned}
\]
Here $\epsilon$ is the empty program.  The production $P ::= C P$
corresponds to the \texttt{Sequence(Vec<BFProg>)} case in my Rust code,
and $C ::= \texttt{[}P\texttt{]}$ corresponds to
\texttt{Loop(Box<BFProg>)}.  In particular, the grammar has no production
for a left or right bracket by itself, just as my Rust enum has no variant
for a standalone bracket.

\item
Prove that \texttt{]><+[]} is not a valid BF program given your formal grammar.
(Note: this is importantly different from why this is not a valid BF program according to your Rust enum!
In the Rust enum case, it should be impossible to construct this program; but in your syntax it should be because
your grammar doesn't allow this string to be part of the grammar.

\paragraph{Answer.}
I will use the following invariant: every string derived from $P$ has
balanced brackets, and every prefix of the string contains at least as many
\texttt{[} symbols as \texttt{]} symbols.

This follows by induction on the grammar.  The empty string has the
property.  Each non-bracket command also has the property.  If $x$ is
derived from $P$, then \texttt{[}$x$\texttt{]} has the property because the
new opening bracket comes before the new closing bracket.  Finally,
concatenating two strings with balanced brackets and the prefix property
preserves both properties.  These cases cover all productions in the
grammar.

The first character of \texttt{]><+[]} is \texttt{]}.  Therefore, its
one-character prefix contains one closing bracket and no opening bracket,
which violates the invariant.  Consequently there is no derivation
$P \Rightarrow^* \texttt{]><+[]}$, so the string is not a valid BF program
according to the grammar.

\item
Describe the semantics for BF programs
by providing a \emph{mathematical} definition of the set of triples \texttt{(INPUT, prog, OUTPUT)} that are valid executions of BF programs.
You will need to start by defining INPUT and OUTPUT.\footnote{\textbf{Hint:} sequence of bytes is a good choice.}
First: do this by assuming you have a deterministic reference interpreter, that is a program that can be used on any INPUT to run the BF program, and produces some output.
You should use this to define a mathematical definition of the semantics.

Second part \textbf{(extra credit)}: Define this without assuming access to a reference interpreter of any kind,
by directly defining the semantics.
Here is a sketch: define a ``configuration'' of a BF program (this should include input, output, and the working tape), define a starting configuration, and then define what it means for a configuration to step to another configuration.
Then state when \texttt{(INPUT, prog, OUTPUT)} is valid by saying something like: there is a sequence of configurations $c_0, c_1, c_2, c_3, \ldots$, where $c_0$ is the starting configuration on input INPUT and program \texttt{prog}, \ldots where you will need to relate each configuration with the previous one.

\paragraph{Answer (using the reference interpreter).}
Let
\[
B = \{0,1,\ldots,255\}
\]
be the set of bytes, and let $B^*$ be the set of finite byte sequences.  An
\texttt{INPUT} and an \texttt{OUTPUT} are both elements of $B^*$.  Let
$\mathcal{P}$ be the set of programs generated by the grammar from part (a).

For $I \in B^*$ and $p \in \mathcal{P}$, write
\[
\operatorname{Run}(p,I) \mathrel{\downarrow} O
\]
when my deterministic \texttt{BFInterpreter}, started with input $I$, an
all-zero tape, data pointer zero, and empty output, terminates normally with
output $O$.  The relation is intentionally partial: a program might run
forever, move left from tape position zero, or request a byte after its input
is exhausted.  In those cases there is no output related to that execution.
The semantics of BF is the following set:
\[
\mathit{Sem}_{\mathrm{BF}}
=
\{(I,p,O) \in B^* \times \mathcal{P} \times B^*
  \mid \operatorname{Run}(p,I) \mathrel{\downarrow} O\}.
\]
Thus a triple is valid exactly when the reference interpreter terminates on
that program and input and produces exactly the stated output bytes.

\paragraph{Answer (direct semantics, extra credit).}
The same relation can be defined without referring to the Rust interpreter.
A configuration is a tuple
\[
\langle K,T,d,I,O\rangle,
\]
where $K$ is a finite list of program nodes still to execute,
$T:\mathbb{N}\to B$ is the tape, $d\in\mathbb{N}$ is the data pointer,
$I\in B^*$ is the unread input, and $O\in B^*$ is the output produced so
far.  I write $q::K$ for a program node followed by the rest of the work
list, $[]$ for the empty list, $T[d\mapsto b]$ for a tape update, and
$O\cdot b$ for appending byte $b$ to the output.

The starting configuration for program $p$ and input $I$ is
\[
c_0 = \langle [p],T_0,0,I,\epsilon\rangle,
\qquad T_0(j)=0\text{ for every }j\in\mathbb{N}.
\]
The structural rules put sequence elements onto the work list and implement
a loop by either discarding it or scheduling its body followed by the loop
again:
\[
\begin{aligned}
\langle \operatorname{Sequence}(q_1,\ldots,q_n)::K,T,d,I,O\rangle
&\to \langle q_1::\cdots::q_n::K,T,d,I,O\rangle,\\
\langle \operatorname{Loop}(q)::K,T,d,I,O\rangle
&\to \langle K,T,d,I,O\rangle &&\text{if }T(d)=0,\\
\langle \operatorname{Loop}(q)::K,T,d,I,O\rangle
&\to \langle q::\operatorname{Loop}(q)::K,T,d,I,O\rangle
&&\text{if }T(d)\ne 0.
\end{aligned}
\]
The primitive-command rules are
\[
\begin{aligned}
\langle \texttt{>}::K,T,d,I,O\rangle
&\to \langle K,T,d+1,I,O\rangle,\\
\langle \texttt{<}::K,T,d,I,O\rangle
&\to \langle K,T,d-1,I,O\rangle &&\text{if }d>0,\\
\langle \texttt{+}::K,T,d,I,O\rangle
&\to \langle K,T[d\mapsto T(d)+1\bmod 256],d,I,O\rangle,\\
\langle \texttt{-}::K,T,d,I,O\rangle
&\to \langle K,T[d\mapsto T(d)-1\bmod 256],d,I,O\rangle,\\
\langle \texttt{.}::K,T,d,I,O\rangle
&\to \langle K,T,d,I,O\cdot T(d)\rangle,\\
\langle \texttt{,}::K,T,d,b::I,O\rangle
&\to \langle K,T[d\mapsto b],d,I,O\rangle.
\end{aligned}
\]
There is no transition for \texttt{<} when $d=0$, or for \texttt{,} when
the unread input is empty; these are precisely the undefined cases used by
my interpreter.  Finally,
\[
(I,p,O)\in\mathit{Sem}_{\mathrm{BF}}
\]
exactly when there are configurations $c_0,c_1,\ldots,c_n$ such that $c_0$
is the starting configuration above, $c_i\to c_{i+1}$ for every $i<n$, and
\[
c_n=\langle [],T,d,I',O\rangle
\]
for some final tape $T$, pointer $d$, and unused input suffix $I'$.  An
infinite transition sequence describes a nonterminating program and therefore
does not produce a valid output triple.

\item
Prove formally that the semantics of BF is deterministic.
You should use your definition above and the definition of deterministic semantics given in class.

\paragraph{Answer.}
We must show that for every $I$, $p$, $O_1$, and $O_2$,
\[
(I,p,O_1)\in\mathit{Sem}_{\mathrm{BF}}
\ \land\
(I,p,O_2)\in\mathit{Sem}_{\mathrm{BF}}
\quad\Longrightarrow\quad
O_1=O_2.
\]
By the definition in part (c), the two premises say
$\operatorname{Run}(p,I)\mathrel{\downarrow}O_1$ and
$\operatorname{Run}(p,I)\mathrel{\downarrow}O_2$.  The reference
interpreter is deterministic: each \texttt{BFProg} variant selects exactly
one branch of \texttt{interpret}; arithmetic and pointer updates have unique
results; sequences execute in their stored order; and a loop's current tape
byte uniquely determines whether the body runs again.  Reading input is also
deterministic because \texttt{,} consumes the unique first unread byte.
Therefore, from the same initial state, every next state is uniquely
determined.

Equivalently, in the direct semantics above, no configuration can match two
different transition rules.  The program node at the head of $K$ determines
the rule, and the two loop premises $T(d)=0$ and $T(d)\ne0$ are mutually
exclusive.  Induction on the number of steps therefore shows that two
executions starting at the same $c_0$ have identical configurations at every
step where both are defined.  If both terminate, their final output components
must be equal.  Hence $O_1=O_2$, which proves that
$\mathit{Sem}_{\mathrm{BF}}$ is deterministic.
\end{enumerate}

% %%%%%%%%%%%%%%%%%%%%%%%%%
% %%%%%   PROBLEM 3   %%%%%
% %%%%%%%%%%%%%%%%%%%%%%%%%

\item
This problem continues the previous two problems.
Something we have not talked about yet in class is parsers.
For this problem, you will write a naive parser for BF programs.
Your parser should be implemented as a function \texttt{parse\_bf}.
It should take a Rust string reference (\texttt{\&str}) that represents a BF program,
and return either \texttt{Some(prog)}, if the program is a valid BF program,
or \texttt{None}, if parsing failed.

Here is a starting skeleton. Please use the signature as follows
(you can modify the body).
\begin{verbatim}
pub fn parse_bf(bf_str: &str) -> Option<BFProg> {
    // define some additional state here ...
    // (what state would be useful to maintain during the parsing?)

    // Iterate through the string. We can use .chars():
    for char in bf_str.chars() {
        // TODO
        unimplemented!()
    }
}
\end{verbatim}

Write at least two unit tests that show your parser seems to be working.

Answer the following in the LaTeX:
is there any program for which displaying it, followed by parsing the resulting display, does NOT result
in the same program that you started with?
Why or why not?

If yes, demonstrate an example of such a program in a unit test.

Your unit tests for this problem should be called \texttt{test\_parser\_1}, \texttt{test\_parser\_2}, \texttt{test\_parser\_3}, etc.

\paragraph{Answer.}
Yes.  In my representation, displaying and then parsing always preserves the
BF command string, but it does not always preserve the exact Rust enum value.
The reason is that \texttt{Display} prints the contents of a
\texttt{Sequence} without printing any marker for the \texttt{Sequence} node
itself.  Consequently, different syntax trees can have the same display.

For example, let
\[
p=\texttt{BFProg::IncByte}.
\]
Displaying $p$ produces \texttt{+}.  My parser represents every complete
input as a top-level sequence, so parsing that string produces
\[
\texttt{BFProg::Sequence(vec![BFProg::IncByte])},
\]
which is not structurally equal to $p$, even though it displays as the same
string and has the same BF behavior.  The unit test \texttt{test\_parser\_3}
demonstrates this example.  The other parser tests check a nested valid
program as well as unmatched brackets and a non-BF character.

% %%%%%%%%%%%%%%%%%%%%%%%%%
% %%%%%   PROBLEM 4   %%%%%
% %%%%%%%%%%%%%%%%%%%%%%%%%

\item
This problem will explore deterministic versus nondeterministic program semantics.

Assume we have two programming languages, \texttt{Lang1} and \texttt{Lang2}.
Assume that the \texttt{INPUT} and the \texttt{OUTPUT} in both languages is a sequence of integers.
Let \texttt{Semantics1} and \texttt{Semantics2} be sets of \texttt{(INPUT, prog, OUTPUT)} pairs, so that
$$
(\texttt{INPUT}, \texttt{prog}, \texttt{OUTPUT}) \in \texttt{Semantics1}
$$
if and only if the semantics oracle for \texttt{Lang1} accepts the program, and similarly for \texttt{Lang2}.

We say that language 2 can \emph{simulate} language 1 if there is a function $f$ that takes any program
$p_1$ in \texttt{Lang1} and produces a program $f(p_1)$ in \texttt{Lang2} such that,
for all \texttt{INPUT} and \texttt{OUTPUT}:
$$
(\texttt{INPUT}, p_1, \texttt{OUTPUT}) \in \texttt{Semantics1}
\text{ iff }
(\texttt{INPUT}, f(p_1), \texttt{OUTPUT}) \in \texttt{Semantics2}.
$$

(\textbf{Note:} iff means ``if and only if''.)

\begin{enumerate}[(a)]
\item Prove that it is impossible for a deterministic programming language to simulate a nondeterministic one in
this sense. That is, prove that if
\texttt{Semantics2} is deterministic, and \texttt{Semantics1} is not, then \texttt{Semantics2} cannot simulate \texttt{Semantics1}.

\paragraph{Answer.}
Assume for contradiction that \texttt{Lang2} does simulate \texttt{Lang1}
through some translation function $f$.  Because \texttt{Semantics1} is not
deterministic, there must be an input $I$, a program $p$, and two different
outputs $O_1\ne O_2$ such that
\[
(I,p,O_1)\in\texttt{Semantics1}
\qquad\text{and}\qquad
(I,p,O_2)\in\texttt{Semantics1}.
\]
In other words, running the same program on the same input can have at least
two possible results.

Now apply the definition of simulation to each of these triples.  It gives
\[
(I,f(p),O_1)\in\texttt{Semantics2}
\qquad\text{and}\qquad
(I,f(p),O_2)\in\texttt{Semantics2}.
\]
The translated program $f(p)$ therefore has two different outputs on the
same input $I$.  This contradicts the assumption that
\texttt{Semantics2} is deterministic.  Thus no such translation $f$ can
exist, and a deterministic language cannot simulate a nondeterministic one
under this definition of simulation.

\item Prove, by giving a specific example, that the other way around is possible:
that is, it is possible for a nondeterministic programming language to simulate a deterministic one.

\paragraph{Answer.}
Consider a small deterministic language \texttt{Lang1} whose programs are
integers.  A program $n$ ignores its input and returns the one-element output
sequence $[n]$.  Its semantics is
\[
\texttt{Semantics1}
=\{(I,n,[n])\mid I\text{ is an integer sequence and }n\in\mathbb{Z}\}.
\]
This language is deterministic because a fixed program $n$ has exactly one
possible output, $[n]$, on any fixed input.

Now define \texttt{Lang2} to have two forms of programs:
\[
\operatorname{Return}(n)
\qquad\text{and}\qquad
\operatorname{Choose}(a,b).
\]
The first form deterministically returns $[n]$.  The second may return either
$[a]$ or $[b]$.  More formally, its semantics contains
\[
\begin{aligned}
(I,\operatorname{Return}(n),[n])&\quad\text{for every }I,n,\\
(I,\operatorname{Choose}(a,b),[a])&\quad\text{for every }I,a,b,\\
(I,\operatorname{Choose}(a,b),[b])&\quad\text{for every }I,a,b.
\end{aligned}
\]
This language is nondeterministic: for example,
$\operatorname{Choose}(0,1)$ can produce either $[0]$ or $[1]$ on the same
input.

Define the translation by
\[
f(n)=\operatorname{Return}(n).
\]
For every input $I$ and output $O$,
\[
(I,n,O)\in\texttt{Semantics1}
\quad\Longleftrightarrow\quad O=[n]
\quad\Longleftrightarrow\quad
(I,\operatorname{Return}(n),O)\in\texttt{Semantics2}.
\]
Therefore \texttt{Lang2} simulates all of \texttt{Lang1}, even though
\texttt{Lang2} also contains other programs that behave nondeterministically.
The simulator does not have to use the nondeterministic feature when it is
translating a deterministic program.

\item
Prove that, despite part (a), it is possible for a deterministic programming language to simulate a nondeterministic
one in the following different sense:
if we are allowed to modify the input by adding additional integers to the \texttt{INPUT} sequence.
The modified criteria for simulation is as follows:
For all $\texttt{INPUT}_1$ and \texttt{OUTPUT},
$$
(\texttt{INPUT}_1, p_1, \texttt{OUTPUT}) \in \texttt{Semantics1}
$$
iff there exists $\texttt{INPUT}_2$ such that:
$$
(\texttt{INPUT}_2, f(p_1), \texttt{OUTPUT}) \in \texttt{Semantics2} \text{ and } \texttt{project}(\texttt{INPUT}_2) = \texttt{INPUT}_1
$$

where above, $\texttt{project}$ denotes an additional function defined by the simulation which takes
$\texttt{INPUT}_2$ and gives back $\texttt{INPUT}_1$.
The intuition should be as follows:
$\texttt{INPUT}_2$ contains some additional input to the deterministic language, beyond what was present in the nondeterministic one.

\paragraph{Answer.}
The extra input can record the choices that a nondeterministic execution
would make.  I will describe the standard construction in which the possible
transitions at each nondeterministic step are numbered by integers.  This
covers the usual finite-choice and countable-choice forms of nondeterminism.

Suppose the original input is
\[
I_1=[x_1,x_2,\ldots,x_n].
\]
An input to the deterministic simulator has the form
\[
I_2=[n,x_1,x_2,\ldots,x_n,r_1,r_2,\ldots,r_k],
\]
where the suffix $R=[r_1,\ldots,r_k]$ is a list of choices.  The first
integer tells us how much of the sequence belongs to the original input.
Define
\[
\operatorname{project}(I_2)=[x_1,x_2,\ldots,x_n]=I_1.
\]
Malformed encodings may be projected arbitrarily; the deterministic
simulator rejects them and therefore produces no output triple from them.

For each nondeterministic program $p$, let $f(p)$ be a deterministic program
that simulates $p$ as follows.  It uses $I_1$ whenever $p$ requests ordinary
input.  Whenever $p$ reaches a nondeterministic choice, $f(p)$ consumes the
next integer $r_j$ from $R$ and uses it to select one of the available
transitions.  For example, if there are $m$ alternatives, it can select the
alternative numbered $r_j\bmod m$.  If the choice sequence ends before the
simulated execution terminates, the simulator rejects that encoded input.
Once the complete input $I_2$ is fixed, every choice is fixed, so $f(p)$ is a
deterministic program.

It remains to prove the required equivalence.

\emph{Forward direction.}
Suppose
\[
(I_1,p,O)\in\texttt{Semantics1}.
\]
Then there is some terminating nondeterministic execution of $p$ on $I_1$
that produces $O$.  Record, in order, the branch selected at every
nondeterministic step of that execution.  Since the execution terminates, it
uses only a finite choice sequence $R$.  Construct $I_2$ by appending that
sequence to the length-prefixed encoding of $I_1$.  The deterministic program
$f(p)$ follows exactly the recorded execution, so
\[
(I_2,f(p),O)\in\texttt{Semantics2}
\qquad\text{and}\qquad
\operatorname{project}(I_2)=I_1.
\]

\emph{Reverse direction.}
Suppose there is an $I_2$ such that
\[
(I_2,f(p),O)\in\texttt{Semantics2}
\qquad\text{and}\qquad
\operatorname{project}(I_2)=I_1.
\]
The extra integers in $I_2$ direct $f(p)$ only along transitions that are
available to the original nondeterministic program.  If we erase the choice
encoding and look at those selected transitions, we obtain a valid execution
of $p$ on $I_1$ with output $O$.  Therefore
\[
(I_1,p,O)\in\texttt{Semantics1}.
\]

Both directions hold, so existentially choosing the additional input recovers
exactly the set of outputs of the nondeterministic program.  The apparent
nondeterminism has not disappeared; it has been moved into the choice suffix
of the deterministic simulator's input.
\end{enumerate}

% %%%%%%%%%%%%%%%%%%%%%%%%%
% %%%%%   PROBLEM 5   %%%%%
% %%%%%%%%%%%%%%%%%%%%%%%%%

\item
This last problem will explore how deep mathematical statements can be embedded in program semantics.

The \emph{Goldbach Conjecture} is an open mathematics problem.
It states that for any even integer greater than $2$, that integer can be expressed
as the sum of two prime numbers.

Write a program in BF (or find one online) that terminates if and only if the Goldbach conjecture is false.

Run your program in your BF interpreter (and check that it seems to not terminate).
Ctrl-C your program after a while and make sure not to include the code which runs this program
in your final submitted code.

Then answer the following questions:

\paragraph{BF program.}
My function \texttt{goldbach\_program()} in \texttt{main.rs} constructs the
required closed BF program.  I used Tom Smeding's Turing-machine simulator
written in BF,
\url{https://git.tomsmeding.com/bfturing/about/}, and specialized its finite
input-loading phase to a fixed 25-state Goldbach machine.  The transition
table is the \texttt{edited} machine in Google DeepMind's formally verified
development,
\url{https://github.com/google-deepmind/cryptids/blob/main/Cryptids/Goldbach21.lean}.
That development proves that this machine, started in state 2 on a blank tape,
halts exactly when some even number greater than 2 has no Goldbach partition.

The Rust construction first runs only the simulator's terminating input loader
on the fixed transition table.  It then converts the resulting tape into BF
initialization commands and appends the BF simulation loop.  This partial
evaluation makes the result a self-contained \texttt{BFProg}: it has no
external input and contains only the eight BF operations.  Constructing the
value does not start the Goldbach search.  I temporarily ran the completed
program through my \texttt{BFInterpreter}; it continued running until I
stopped it with Ctrl-C.  The temporary call that launched it was then removed,
so neither \texttt{main()} nor the submitted tests start the nonterminating
search.

The reason for the termination equivalence is a composition of two facts.  The
BF simulator halts exactly when the encoded Turing machine halts, and the
encoded machine halts exactly when it finds a Goldbach counterexample.
Therefore the constructed BF program halts if and only if the Goldbach
conjecture is false.

\begin{enumerate}[(a)]
\item Do you think that this is anything special about the Goldbach conjecture? Or can all mathematical statements
be expressed as whether a particular program terminates?

\paragraph{Answer.}
There is nothing uniquely special about Goldbach's conjecture here.  Its
useful feature is that a counterexample can be recognized by a finite
computation.  For a proposed even number $n$, a program can check every pair
$p+(n-p)=n$ and decide whether either number is prime.  A program can therefore
test $4,6,8,\ldots$ in order and halt when it finds the first counterexample.
If no counterexample exists, that search runs forever.

The same construction works for any statement of the form
\[
\forall n\in\mathbb{N},\;P(n),
\]
provided that $P(n)$ is decidable for each particular $n$.  The program simply
enumerates $n$ and halts when $P(n)$ is false.  Many number-theoretic
conjectures have this form, so Goldbach is an especially understandable
example rather than an exceptional one.

I would not claim, without qualification, that every possible mathematical
statement is directly equivalent to one ordinary halting question by this
same counterexample-search construction.  Statements with more complicated
alternations of ``for every'' and ``there exists'' can lie higher in the
arithmetical hierarchy.  Also, searching through formal proofs makes a program
halt when a proof is found, but an independent true statement might have no
proof in the chosen formal system.  The accurate conclusion is that halting
already expresses every computably searchable existence claim, including the
existence of a Goldbach counterexample, and that class is broad enough to
contain extremely difficult and even undecidable families of questions.

\item Give your best answer to the following:
why is this not a problem for program semantics?

That is, why is it not an issue that program semantics can express arbitrarily complex mathematical statements, that may
be very difficult or even impossible to resolve? \emph{(Note: we know there are some mathematical statements that cannot ever
be proven one way or the other, although it is unknown if the Goldbach conjecture is one of such statements, it would seem to be plausible.)}
Doesn't this make program semantics totally intractable as a subject?
Yet, we can define program semantics for practical languages like C, Python, or Rust.
Explain.

\paragraph{Answer.}
This is not a problem because defining a program's semantics is different from
having an algorithm that answers every question about that program.  The
semantics only needs to state precisely how one configuration steps to the
next and what counts as an observable result.  Each individual BF step is
simple: move the pointer, change a byte, read or write a byte, or choose whether
to repeat a loop.  The difficult behavior appears when those simple rules are
composed for an unbounded number of steps.

For the Goldbach program, the semantics unambiguously tells us what every
finite execution prefix looks like.  It also gives a precise conditional
fact: if a counterexample is reached, the program terminates; if every even
number has a prime partition, the execution is infinite.  The semantics does
not need to decide which mathematical alternative is true in order to define
the program's behavior correctly.

The same distinction applies to practical languages such as C, Python, and
Rust.  We can define evaluation, memory, function calls, exceptions, and I/O
precisely, even though no tool can decide termination or every other semantic
property for every program in those languages.  Static analyzers, type
systems, proof assistants, and human proofs remain useful by being sound on
the cases they can handle; they do not have to be complete decision procedures
for all programs.

In fact, the possibility of embedding hard mathematics is a consequence of a
language being computationally expressive.  Program semantics gives us a
rigorous vocabulary for reasoning about that expressiveness.  It does not make
all reasoning easy, but it lets us prove local rules, equivalences, safety
properties, and correctness results where suitable invariants are available.
So undecidability limits what can be automated universally; it does not prevent
the language from having a clear and useful semantics.
\end{enumerate}

% %%%%%%%%%%%%%%%%%%%%%%%%%
% %%%%%%%%%%%%%%%%%%%%%%%%%

\end{enumerate}

\newpage

\noindent
\textbf{Submission instructions:}
\begin{itemize}
\item Upload your solutions (as a PDF) and your code (as a \texttt{.zip} file, including the entire repository directory including config files like \texttt{Cargo.toml}, as with HW0) in Gradescope.
\item Your PDF should be a single file \texttt{hw1.pdf}. If you use this LaTeX template to create your solutions, please remove the problem statements and include only your solutions.
\item Your code (functions and types as described above) should be contained in \texttt{main.rs}.
(You may also add utility modules, like \texttt{mod util;} then \texttt{use util::...}, if you feel that you need them; but you should put all the main functions with the names as above in \texttt{main.rs}.) Please double check on Piazza before importing any external crates (in your \texttt{Cargo.toml}).

\item The \texttt{main.rs} should contain your solutions for parts 1, 3, and 5.

\item Please include all of the function names and signatures as listed in the document above, and do not modify any function signatures. You are welcome to add additional functions and tests -- but if you do so, please include all additional tests only at the bottom of the file.

\item Make sure that \texttt{cargo build}, \texttt{cargo clippy}, \texttt{cargo test}, and \texttt{cargo fmt} (or \texttt{cargo fmt --check}) pass without any errors.
\end{itemize}

\end{document}
